% Original successor argument: audit of LM6 A.50--A.52 and Lead 8.
\subsection{The tetrahedral curve and the arithmetic part of Lead 8}

The group denoted by $\mathcal K$ in this subsection is the group in
the question, identified above with
$C_{\widehat\Gamma_{0,[12]}}(A_5)$; in particular it is pure.
The letter $K$ without calligraphic font retains its meaning as the
level-five free subgroup of $\mathrm{PSL}_2(\mathbb Z)$.
All completions in this subsection are full profinite completions.

\statement{B.29}{Proved}{The precise arithmetic normalizer restriction}
Let $X/k$ be a smooth affine hyperbolic curve over a number field, with
geometric fundamental group $F=\pi_1(X_{\overline k})$ and arithmetic
fundamental group $E=\pi_1(X)$.  Let $1\ne A\leq F$ be closed and
procyclic, and suppose that a closed subgroup $N\leq E$ normalizes $A$
and has image $G_L\leq G_k$, with $L/k$ finite.  For a finite Galois
extension $M/\mathbb Q$ containing $L$, let $\mathcal P_M$ be the
rational primes unramified in $M$ whose Frobenius class is the class
of a complex conjugation in $\operatorname{Gal}(M/\mathbb Q)$.
Then either $A$ is contained in a unique cusp inertia group, or
\[
                         A_\ell=1\quad(\ell\in\mathcal P_M).
                                                        \tag{B.29.1}
\]
Consequently cofinite prime support implies cuspidality.  If $L$ is
totally real, every such $A$ is cuspidal, without a prime-support
condition.  These are cuspidality conclusions, not nonexistence
conclusions in an arbitrary curve group.

\begin{proof}
Here is an independent proof of the three unrefereed rows LM6
A.50--A.52 used in this claim, with one missing preliminary step
made explicit.  Replace $N$ by the closed subgroup $AN$ first.
Its arithmetic image is unchanged, and its action on $A$ extends
the original action because $A$ is abelian.  Thus we may and do
have $A\leq D:=N\cap F$.  Without this replacement the Sylow subgroup
containing $A_\ell$ in the printed proof of LM6 A.50 need not exist.

The affine geometric group $F$ is free profinite and torsion-free.
For a prime with $A_\ell\ne1$, therefore, $A_\ell\simeq\mathbb Z_\ell$.
Write $\theta_\ell:N\to\mathbb Z_\ell^\times$ for conjugation on this
subgroup.  In a pro-$\ell$ Sylow subgroup $S$ of $D$ containing
$A_\ell$, each group $\langle A_\ell,g\rangle$, $g\in S$, is both
free pro-$\ell$ and metabelian.  It is procyclic: a free pro-$\ell$
group of rank at least two has nonmetabelian finite $\ell$-group
quotients.  In particular $g$ centralizes $A_\ell$.  The image of
a Sylow subgroup is a Sylow subgroup of the image, so
$\theta_\ell(D)$ has trivial pro-$\ell$ part.  Put $d=\ell-1$ for
odd $\ell$ and $d=1$ for $\ell=2$.  There is consequently a continuous
character on the actual compact quotient $N/D=G_L$,
\[
 \lambda_\ell:G_L\longrightarrow\mathbb Z_\ell^\times,
       \qquad \lambda_\ell\circ p=\theta_\ell^d.       \tag{B.29.2}
\]
No section of $N\to G_L$ is required here.

Choose an open characteristic subgroup $V$ of $F$ not containing
a generator of $A_\ell$, and set $H=VA_\ell$.  The finite cyclic
quotient $H/V$ shows that $A_\ell\to H^{\rm ab}$ is nonzero.
Since $H^{\rm ab}$ is torsion-free, this map is injective on
$A_\ell$.  Descend this finite geometric cover to a finite
extension and realize its arithmetic fundamental group as an
open subgroup $E_H\leq E$ with $E_H\cap F=H$.  The image of
$N\cap E_H$ in $G_L$ is open; write it as $G_{L_1}$.
Conjugation on $H^{\rm ab}$ by $N\cap E_H$ factors through
$G_{L_1}$ because its geometric kernel is contained in $H$.
Thus the nonzero $\mathbb Q_\ell$-line furnished by $A_\ell$
is a genuine arithmetic line with character $\psi$, and
\[
                   \lambda_\ell|_{G_{L_1}}=\psi^d.          \tag{B.29.3}
\]

The kernel of geometric homology onto the Tate module of the
Jacobian of the smooth compactification of this cover is the
cuspidal permutation module, modulo its single degree relation,
tensored with $\mathbb Z_\ell(1)$.  If our line has zero image in
the Jacobian Tate module, its character is cyclotomic after
splitting the finitely many cusp fields.  Stix's essentially
cyclotomic normalizer criterion, Theorem~20 of
\cite{CEStixCuspidal2008}, puts $A_\ell$ inside a cusp inertia
group $I$ of $X$.  The abelian group $A$ then lies in $I$ by
Proposition~11 and Corollaries~12--13 of that paper.  These are
the inertia-intersection and normalizer statements required
here; its Corollary~7, cited for that purpose in LM6 A.14,
instead concerns the Kummer map and is not the required reference.

In the other case the image of our line in the Jacobian Tate
module is nonzero and injective over $\mathbb Q_\ell$.
The character $\psi$ has Frobenius weight one and Hodge--Tate
weights in $\{0,1\}$, with cyclotomic weight one in the latter
convention.  It is unramified away from finitely many places.
The same finiteness and Hodge--Tate properties hold for
$\lambda_\ell$; Hodge--Tate descent through a finite extension
is the descent statement on page III--32 of
\cite{CESerreAbelian1989}.

For completeness, the one-dimensional character calculation
used next does not require that Frobenius eigenvalues be rational.
Let $\lambda:G_M\to\mathbb Q_\ell^\times$ be Hodge--Tate.
Local algebraicity, followed by global reciprocity, supplies a
$\mathbb Q_\ell$-defined algebraic character of
$(\operatorname{Res}_{M/\mathbb Q}\mathbb G_m)_{\mathbb Q_\ell}$
with integer coefficients $n_\tau$, indexed by embeddings of
$M$ in an algebraic closure.  It is trivial on a subgroup of
finite index of the global unit group.  Taking complex absolute
values and applying the logarithmic form of Dirichlet's unit
theorem gives
\[
             n_\tau+n_{c\tau}=2a\quad\hbox{for every }\tau. \tag{B.29.4}
\]
If $M$ is totally real, it gives directly $n_\tau=a$ for every
$\tau$.  Otherwise, when $\ell\in\mathcal P_M$, choose the place
above $\ell$ whose Frobenius is the chosen $c$.  The coefficient
vector is fixed by its decomposition group, because the
algebraic character is defined over $\mathbb Q_\ell$.
Hence $n_{c\tau}=n_\tau$, and (B.29.4) again gives one integer
$n_\tau=a$ for every $\tau$.  Removing $\chi_\ell^a$ leaves a
character trivial on an open subgroup of every local unit
group; global reciprocity makes it a finite ray-class character.
Thus
\[
                         \lambda=\varepsilon\chi_\ell^a,
                       \qquad |\varepsilon(G_M)|<\infty.  \tag{B.29.5}
\]
The input to this calculation is Tate's local algebraicity
theorem and the modulus construction, Chapter~III,
pages III--7, III--9 and III--11 of
\cite{CESerreAbelian1989}; the unit-character lattice is
also computed there in Chapter~II, pages II--29--30.

Apply (B.29.5) to (B.29.2), restricted to $G_M$, and pass
to a common finite extension killing $\varepsilon$ and on which
(B.29.3) holds.  The equality $\psi^d=\chi_\ell^a$ has two
incompatible consequences.  Frobenius weights at a good
place away from $\ell$ give $d=2a$, while Hodge--Tate weights
give $a\in\{0,d\}$.  This contradicts $d>0$.  Thus a
noncuspidal $A$ has $A_\ell=1$ at the primes in (B.29.1).
Chebotarev gives that set density
$|\operatorname{Cl}(c)|/[M:\mathbb Q]>0$, proving the
cofinite-support assertion.  If $L$ is totally real, take its
totally real Galois closure.  The same character calculation
works at every $\ell$, so any nonzero prime component gives
the contradiction.  Finally distinct cusp inertia groups
have trivial intersection by Stix, Proposition~11, which
proves uniqueness.
\end{proof}

\statement{B.30}{Proved}{The tetrahedral centralizer has its full profinite topology}
Let $T=A_4$ act on the twelve vertices and put
$F_T=C_{\Gamma_{0,[12]}}(T)$.  The natural map
\[
              \widehat{F_T}\longrightarrow\widehat\Gamma_{0,[12]}
                                                               \tag{B.30.1}
\]
is injective, and its image is exactly
$C_{\widehat\Gamma_{0,[12]}}(T)$.  In particular $F_T\simeq F_2$
gives $\overline{F_T}\simeq\widehat F_2$, and the pure subgroup
$F_T\cap\Gamma_{0,12}\simeq F_{13}$ also has its full profinite
topology induced from the ambient mapping class group.

\begin{proof}
We separate the topology assertion from the density assertion;
neither follows merely by replacing a subgroup with its closure.
The induction below rechecks the argument of LM4 F.7.4--F.7.6;
it is not claimed as a new topology argument.
First recall the valid prime-power specialization of the
symmetric congruence argument in \cite{CEBoggiSymmetric2026}.
For a finite $p$-subgroup $J$ of a genus-zero mapping class group,
its centralizer has the full profinite topology induced by
the ambient congruence topology.  Here is why the problematic
finite-intersection step is not needed for this specialization.
Remove the ramification points of the $J$-action.  In the
resulting free action, restriction to the normal finite-index
surface subgroup proves injectivity of the completed
centralizer representation: an automorphism of a nonabelian
profinite surface group restricting to the identity on an open
normal subgroup is the identity.  The finite-index quotient
mapping class group is a genus-zero group, for which CSP holds.
Filling the ramification points has configuration-group kernel
$\widehat{\Pi[r]}$, by the profinite Birman exact sequence.
The subgroup $\Pi[r]\rtimes J$ is residually finite, good and
of finite virtual cohomological type: its finite-index ordered
configuration group is an iterated extension of finitely
generated free groups.  Discretely its centralizer of $J$ is
$Z(J)$, by the ramification-filling comparison.  Apply
Boggi--Zalesskii, Theorem~A(iii), \emph{to the one finite
$p$-subgroup $J$}, to get the same centralizer in its full
completion.  Hence no element of $\widehat{\Pi[r]}$ is lost
on the completed symmetric subgroup.  This proves the required
injectivity and full topology.  These are precisely the free
and ramified arguments of Sections~3.1--3.2 of
\cite{CEBoggiSymmetric2026}, with Proposition~2.4 replaced
only by its valid $p$-subgroup case.  Goodness and the
configuration Birman comparison used there are also treated
in \cite{CEBoggiCongruence2020}.  The ambient genus-zero
congruence completion equals its full completion throughout.

There are two elementary topology transfers.  If a subgroup
$C$ has the full induced topology, a finite extension $N$
of $C$ has it as soon as $C$ is open for the induced topology
on $N$.  Indeed each prescribed finite-index subgroup of $N$
contains a finite-index subgroup of $C$, and intersecting
the corresponding ambient neighborhoods gives the required
neighborhood in $N$.  For $C=C_\Gamma(J)$ and
$N=N_\Gamma(J)$, openness follows by choosing an ambient
normal finite-index $U$ with $U\cap J=1$: an element of
$N\cap U$ centralizes $J$.  Second, for a finite normal
subgroup $J\trianglelefteq N$ in such an embedded group,
the topology on $N/J$ is its full topology.  For a subgroup
$D\leq N/J$ with full induced topology, its inverse image
has full topology too.  To check the last assertion without
assuming it, choose an ambient open subgroup avoiding $J$.
On its intersection with $N$, the quotient map is an
isomorphism onto a finite-index subgroup of $N/J$, with
the induced topologies; apply the first transfer to the
remaining finite extension.

Apply the $p$-subgroup result to $V=V_4\trianglelefteq T$.
The group $N=N_\Gamma(V)$ has full induced topology.
Its quotient $H=N/V$ is the finite-index subgroup of
$\Gamma_{0,[6]}$ preserving the $V$-cover: the quotient
has three images of the marked free orbits and three
branch points.  This is the discrete normalizer exact
sequence, so the induced full topology on $H$ is also
the topology obtained from $\Gamma_{0,[6]}$.
The subgroup $T/V\simeq C_3$ lies in $H$.
Apply the same valid prime-power argument to this $C_3$
in $\Gamma_{0,[6]}$.  Its normalizer, and hence its
finite-index intersection $N_H(C_3)$, has full induced
topology.  Pulling back through the finite kernel $V$
gives full topology on
\[
     N_\Gamma(T)=\{n\in N_\Gamma(V):nV\in N_H(T/V)\}.
\]
Taking the finite-index centralizer proves injectivity in
(B.30.1).  This argument uses two prime-power cases, rather
than assuming the all-solvable normalizer-topology hypothesis
which it is intended to verify.

For density, Theorem~A(iii) of
\cite{CEBoggiZalesskii2025}, first for $V$ in $\Gamma$,
identifies $N_{\widehat\Gamma}(V)$ with $\widehat N$;
the identification with the full completion is justified
by the topology just proved.  After quotienting by $V$,
apply the same theorem to $C_3$ in $H$.  The group $H$
is good, residually finite and of finite virtual
cohomological type, as a finite-index genus-zero mapping
class group.  Pullback therefore gives
$N_{\widehat\Gamma}(T)=\overline{N_\Gamma(T)}$.
The kernel of conjugation on the finite group $T$ is open
in each normalizer.  Density in this open kernel yields
$C_{\widehat\Gamma}(T)=\overline{C_\Gamma(T)}$.

Finally the discrete centralizer is the point-pushing
group of one ordinary orbit in the tetrahedral quotient
sphere with its three branch values removed, hence $F_2$.
Equivalently the discrete centralizer exact sequence has
kernel $Z(T)=1$ and image the pure four-puncture sphere
group.  Its puncture-permutation image is the regular
group $T$, of order twelve.  Intersecting the ambient
open pure subgroup therefore gives an index-twelve
free subgroup, of rank $1+12(2-1)=13$; both its embedding
and its full induced topology follow by restriction.
\end{proof}

\statement{B.31}{Proved}{A totally real arithmetic model at the icosahedral point}
For the tetrahedral action generated by
\[
       u\longmapsto-u,\qquad u\longmapsto1/u,
             \qquad u\longmapsto\frac{u+i}{u-i},
\]
the curve of unlabelled free orbits is the $\mathbb Q$-curve
\[
 C_T=\mathbb P^1_v\setminus\{\infty,6\sqrt{-3},-6\sqrt{-3}\},
 \qquad
 v=\frac{u^{12}-33u^8-33u^4+1}{u^2(u^4-1)^2}.              \tag{B.31.1}
\]
In the embedding used for the cross-ratio computation, the
icosahedral point is
$v_0=-22\sqrt5/5\in C_T(\mathbb Q(\sqrt5))$.
The injection
\[
 \pi_1(C_{T,\overline{\mathbb Q}},\overline v_0)
       \simeq\widehat F_2\longrightarrow\widehat\Gamma_{0,[12]}
                                                               \tag{B.31.2}
\]
has image $\overline{F_T}$.  For the actual Galois action
defined by the rational point $v_0$ and its image family,
$\mathcal K\leq\overline{F_T}$ is stable under the full group
$G_{\mathbb Q(\sqrt5)}$.

\begin{proof}
Put $f=u^{12}-33u^8-33u^4+1$, $g=u^2(u^4-1)^2$ and
$h=u^8+14u^4+1$.  Direct substitution in the three
generators gives invariance of $f/g$; degree twelve
then identifies its function field with the invariant
field.  The exact identities
\[
 h^3-f^2=108g^2,\qquad
       \frac{d}{du}(f/g)=\frac{2h^2}{u^3(u^4-1)^3}         \tag{B.31.3}
\]
give ramification index two above $\infty$ and index
three above each of $\pm6\sqrt{-3}$.  These are all
the exceptional orbits.  In particular the divisor
\[
                               f(u)-vg(u)=0               \tag{B.31.4}
\]
is a family of twelve distinct unordered points on
$\mathbb P^1$ over the displayed $\mathbb Q$-curve.

For an exact basepoint computation let
$q=e^{\pi i/10}$, so $q^8-q^6+q^4-q^2+1=0$, and
$u_0=i(2\cos(\pi/10)+2\cos(\pi/5))$.
Reduction in $\mathbb Q[q]/(\Phi_{20})$ gives
\[
 u_0=q^7+q^6+q^4+q^3,\qquad
 v(u_0)=\frac{44q^6-44q^4-22}{5},\qquad
                       v(u_0)^2=\frac{484}{5}.
\]
The chosen embedding gives the negative real root,
$v(u_0)=-22\sqrt5/5$.  The standalone exact calculation
\texttt{work\_CE/code/check\_tetra\_field.py} verifies
the group order, all twelve invariance identities,
(B.31.3), and these reductions; it uses no floating
point equality.

The analytic monodromy of (B.31.4) is the discrete
point-pushing description of $F_T$ in B.30.  Comparison
of complex and algebraic fundamental groups followed
by (B.30.1) proves (B.31.2), including injectivity of
the full completion.  This avoids identifying a
closure with an abstract free profinite group without
first proving its induced topology.

It remains to check the arithmetic subgroup, since
an unspecified open subgroup of a totally real
Galois group need not correspond to a totally real
field.  The tetrahedral subgroup scheme is stable
over $\mathbb Q$: complex conjugation fixes the first
two generators and sends the third to its composite
with $u\mapsto1/u$; all its coefficients lie in
$\mathbb Q(i)$.  At $v_0$, the entire geometric
automorphism group of the twelve-point divisor is
$A_5$, so it too is stable under $G_{\mathbb Q(\sqrt5)}$.
The rational point $v_0$ gives a section of the
arithmetic fundamental group of $C_T$.  Applying the
family morphism to the unordered moduli stack gives
a section of its arithmetic fundamental group, whose
conjugation action on $\widehat\Gamma_{0,[12]}$
preserves both $T$ and $A_5$.  It also preserves the
pure subgroup, which is the kernel of the finite
permutation representation.  Thus it preserves
$\overline{F_T}=C_{\widehat\Gamma}(T)$ and
$\mathcal K=C_{\widehat\Gamma}(A_5)$, and its action
on the former is exactly the action in (B.31.2).
Changing the chosen geometric basepoint conjugates
these based maps simultaneously and changes none
of the subgroup assertions.
\end{proof}

\statement{B.32}{Proved}{No tetrahedral cusp inertia meets the icosahedral kernel}
For every cusp inertia group $I$ of
$\pi_1(C_{T,\overline{\mathbb Q}},\overline v_0)$,
including every conjugate inertia group,
\[
                             I\cap\mathcal K=1.           \tag{B.32.1}
\]
The statement concerns the whole $\widehat{\mathbb Z}$
inertia, including subgroups supported at one prime.

\begin{proof}
This gives an independent proof of the cusp-free conclusion
used in LM4 F.8; the conclusion was already asserted there.
Fix a subgroup $C_5\leq A_5$ and let $V$ be its two
fixed vertices together with one of its five-element
orbits.  Thus $|V|=7$.  The quotient of this punctured
sphere by $C_5$ has three punctures: the two poles
and the free orbit.  Its discrete centralizer of
$C_5$ is finite by the normalizer exact sequence
over that rigid quotient.  Theorem~A(iii) of
\cite{CEBoggiZalesskii2025}, applied to this one
$5$-subgroup, makes its full profinite centralizer
the same finite group.  The full completion of
the pure seven-puncture sphere group is torsion-free:
goodness and finite cohomological dimension of its
iterated free-group configuration tower give finite
continuous cohomological dimension at each prime.
Consequently the forgetful map
\[
 \operatorname{forg}_V:\widehat\Gamma_{0,12}
                         \longrightarrow\widehat\Gamma_{0,7}
\]
sends every $\kappa\in\mathcal K$ to $1$.  This uses
only purity, equivariance for $C_5$, and the prime-power
comparison, not a centralizer comparison for $A_5$.

Let $\eta$ be a positive meridian of a cusp of $C_T$.
At the cusp $\infty$, the twelve points coalesce in
six pairs; at either finite cusp they coalesce in
four triples.  Write $m=2$ or $m=3$, respectively.
The monodromy $\eta^m$ is pure and is the product
of the positive Dehn twists about the disjoint
cluster boundaries, with exponent one on each.
Indeed in a local coordinate at a ramification
point the quotient map has form $z\mapsto z^m$.
After the base change of degree $m$, one meridian
rotates that cluster once and hence makes its
boundary full twist.  The disjoint local disks
give the stated product.

For \emph{every} seven-element subset $W$ of the
twelve vertices, one cluster contains at least two
members of $W$: there are at most six clusters.
Moreover $W$ meets at least three clusters, since
two clusters contain at most six vertices.  Choose
two members in the first cluster and one in each
of two other clusters.  Forgetting to these four
points sends $\eta^m$ to exactly one essential
Dehn twist on a four-puncture sphere; all the
other cluster twists become peripheral or trivial.
That Dehn twist generates a full procyclic
$\widehat{\mathbb Z}$ in the completed pure
four-puncture sphere group $\widehat F_2$.
For example it is a primitive cusp generator of
the latter free group, whose cyclic subgroup is
a retract before and after completion.  Therefore
the image of $\eta^{ma}$ is nontrivial for every
$0\ne a\in\widehat{\mathbb Z}$.

Now suppose $\kappa=\gamma\eta^a\gamma^{-1}$ belongs
to $\mathcal K$, where $\gamma\in\widehat F_T$.
Choose $\gamma_0\in F_T$ with the same puncture
permutation as $\gamma$, and write $\gamma=\gamma_0p$
with $p$ pure.  Set $W=\gamma_0^{-1}V$.  The naturality
of forgetting for the discrete mapping class
$\gamma_0$ gives an isomorphism between the two
seven-puncture targets under which
$\operatorname{forg}_V(\kappa^m)$ is conjugate to
$\operatorname{forg}_W(\eta^{ma})$; the conjugating
element is the defined pure image of $p$.
The former equals $1$ by the first paragraph.
The four-point test in the preceding paragraph
therefore gives $a=0$.  This proves (B.32.1).

The use of seven points matters.  A five-orbit
test or the abelianization alone can miss a
partition into six antipodal pairs.  The cardinality
argument above works for every conjugate cluster
partition, so neither that exceptional partition
nor a prime-supported exponent escapes the proof.
\end{proof}

\statement{B.33}{Proved}{The abelian case of Lead 8 and the cofinite-support exclusion}
For the actual subgroup $\mathcal K$ of the question,
\[
                         \mathcal K\text{ abelian}
                           \quad\Longrightarrow\quad\mathcal K=1.
                                                               \tag{B.33.1}
\]
Moreover $\mathcal K$ has no nontrivial closed procyclic
subgroup with open arithmetic normalizer on $C_T$
and nonzero pro-$p$ component at all but finitely
many rational primes.  The same exclusion holds
with arbitrary prime support if the normalizer
projects onto the Galois group of a totally real
number field.

\begin{proof}
By B.30--B.31, $\mathcal K$ is a closed subgroup
of the full free profinite group
$\pi_1(C_{T,\overline{\mathbb Q}})$.
If it is abelian, each of its Sylow subgroups
is an abelian free pro-$p$ group, hence is
$1$ or $\mathbb Z_p$.  An abelian profinite
group is the product of its Sylow subgroups,
so
\[
                        \mathcal K\simeq\prod_{p\in S}\mathbb Z_p
\]
is procyclic.  B.31 supplies its normalizer
$\mathcal K\rtimes G_{\mathbb Q(\sqrt5)}$
inside the arithmetic fundamental group of
$C_T$.  Since $\mathbb Q(\sqrt5)$ is totally
real, B.29 makes any nontrivial such group
cuspidal.  This contradicts B.32 and proves
(B.33.1).

For a general nontrivial procyclic subgroup
$A\leq\mathcal K$ with the stated open
arithmetic normalizer, apply the cofinite-support
part of B.29 and then B.32.  If its normalizer
field is totally real, use the last part of
B.29 instead.  Thus no assumption that an
arbitrary open subgroup of $G_{\mathbb Q(\sqrt5)}$
is itself totally real has been made.

The character calculation is also explicit:
in the hypothetical noncuspidal abelian case,
at any nonzero prime component its normalizer
character has a power $\lambda_\ell$ as in
(B.29.2).  Total reality forces
$\lambda_\ell=\varepsilon\chi_\ell^a$ after
restriction to a finite Galois extension,
whereas finite-cover Jacobian homology would
force both $2a=d$ and $a\in\{0,d\}$.  Hence
there is no such nonzero character.  Merely
calling the character ``not cyclotomic'' or
``not CM'' would not supply this contradiction.
\end{proof}

\statement{B.34}{Open}{The missing commutativity assertion in Lead 8}
It remains undecided here whether $\mathcal K$ is
abelian.  No two noncommuting elements of
$\mathcal K$ have been constructed either.
B.33 proves that this is equivalent, for this
particular group, to the unresolved assertion
$\mathcal K=1$.

The arithmetic restriction supplies no way to
extract from an arbitrary nonabelian
$G_{\mathbb Q(\sqrt5)}$-stable closed subgroup
of $\widehat F_2$ a nontrivial procyclic subgroup
with open arithmetic normalizer.  A chosen
element's procyclic closure need not have an
open stabilizer.  Similarly, projecting
$\mathcal K$ to a pro-$p$ or metabelian
quotient does not establish commutativity
in the full group.  These are the exact
remaining steps in the two arithmetic
approaches attempted here.  The finite
cross-ratio and rigid-subconfiguration tests
in the other subsections impose necessary
constraints but do not repair this inference.
